\structural{ВИСНОВКИ}

У роботі побудовано імітаційну модель системи посадки літаків в аеропорту на основі формального опису стохастичною мережею Петрі та визначено кількість посадкових смуг, за якої час окупності додаткових смуг мінімальний.

Концептуальна модель виділяє потік літаків, посадкові смуги, чергу очікування і фінансову підсистему. Формальний опис -- мережа Петрі з 13~позицій і 10~переходів. Головна складність задачі -- індивідуальні таймери дозаправки й відльоту кожного літака. У звичайній мережі Петрі маркери безіменні, і таймер одного літака забирав би з черги інший. Тому мережу доповнено маркерами-об'єктами та дугами, що беруть маркери одного об'єкта. Це невелике розширення формалізму зберегло модель компактною: кожен таймер -- це перехід із затримкою і пара переходів, що перевіряють, чи літак досі в черзі.

Алгоритм імітації реалізовано мовою TypeScript без бібліотек імітаційного моделювання. Він просуває час від події до події, підтримує багатоканальні переходи, інформаційні дуги та розв'язання конфліктів за пріоритетом і ймовірністю. Коректність алгоритму підтверджено модульними тестами, зокрема відтворенням теоретичних характеристик СМО M/M/1, а коректність моделі -- верифікацією на восьми наборах вхідних даних, де результати узгоджуються з формулою Літтла і граничними випадками.

Експериментальне дослідження проведено з прогонами тривалістю рік модельного часу після доби розгону, по 51~прогону для кожної кількості смуг, що забезпечує точність річного прибутку 1~млн~грн з довірчою ймовірністю 0,95. Результати такі:
\begin{itemize}
\item за двох смуг аеропорт збитковий ($-42{,}5$~млн~грн/рік): смуги перевантажені, майже кожен літак дозаправляється, а дві третини відлітають;
\item час окупності мінімальний за семи смуг, тобто за п'яти додаткових: 0,93~міс. за приростом прибутку і 1,19~міс. за загальним прибутком; витрати 15~млн~грн повертаються менш ніж за місяць;
\item за семи смуг жоден літак не відлітає, середнє очікування скорочується з 132 до 9~хв, а річний прибуток становить 151~млн~грн;
\item пробний запуск показав, що оптимум залежить від потоку: за 8~літаків на годину найкраще 5~смуг, за 12 -- 8~смуг; тому за очікуваного зростання потоку варто передбачити шосту додаткову смугу.
\end{itemize}

Прийняті припущення: пороги та грошові величини вигляду «$a \pm b$» розподілені рівномірно, літак дозаправляється не більше одного разу, черга обслуговується в порядку прибуття, у прибуток входять лише дохід від посадок і витрати на дозаправку. Не відтворено інших витрат аеропорту (утримання смуг, персонал), бо умова їх не задає; за їх наявності оптимальна кількість смуг може зменшитись.

Імітаційне моделювання виявилось доречним: воно врахувало нетерплячість літаків з двома порогами, випадкові доходи й витрати та дало оцінку точності результату, чого не дають формули теорії масового обслуговування. Моделювання вважаю успішним. Подальше дослідження варто спрямувати на зміну інтенсивності потоку протягом доби та на правила черги, наприклад пріоритет посадки літакам, що наближаються до порогу дозаправки.
